\documentclass[a4paper,10pt]{article}
\newcommand\subdir{/home/koen/LaTeX-setup/}
\input{\subdir/LaTeX-files/imports.tex}
\input{\subdir/LaTeX-files/master-internship.tex}
\usepackage{\subdir/LaTeX-files/aastex}
\usepackage{natbib}
\bibliographystyle{\subdir/LaTeX-files/astroadstitle.bst}

%Set the title, Author and date
\title{Notes Week 32}
\author{Koen Essers}
\tutorial{Master Internship}
\date{\today}

%Set the header style
\header{\tutorialstring}

%Begin the document
\begin{document}
\setuplayout
\section{Grid}

\begin{figure}[ht]
	\centering
	\incplot{w32-R-M-s-q-panels}
\end{figure}

\section{Mean Molecular Weight}

\subsection{Giant Branch Dip}

Using the \lstinline[style=output, breaklines=true]|EOS_plotter| routine from \emph{MESA}, it
is possible to clearly plot the different EOS tables that MESA uses, and in which regions in
the \(\rho\)-\(T\) diagram.
Below, I show this, and I plot the \(\rho\)-\(T\) profile for
an evolved RGB star over these tables.
\begin{figure}[ht]
	\marginfignote{0}{The RGB profile is colored, where very red regions correspond to regions with relatively large mean molecular weight \(\mu \).
		Blue regions correspond to the lowest mean molecular weights, and white corresponds to an average mean molecular weight of the envelope.
		I think this figure quite clearly shows why we see a pit in the mean molecular weight.
		The core boundary happens inside the Skye EoS table \citep{jermyn2021}.
		This table \emph{assumes} complete ionization.
		Once the profile crosses over to the boundary region, we see that the mean molecular weight increases until it crosses
		the \href{https://freeeos.sourceforge.net}{FreeEOS} regime, which does \emph{not assume} complete ionization.
	}
	\centering
	\incplot{w32-eos}
\end{figure}

\bibliography{master.bib}
\end{document}
